\documentclass[12pt]{article}

% -------------------- Idioma y codificación --------------------
\usepackage[spanish]{babel}
\addto\captionsspanish{\renewcommand{\tablename}{Tabla}}
\usepackage[utf8]{inputenc}
\usepackage[T1]{fontenc}

% -------------------- Formato de página --------------------
\usepackage{geometry}
\geometry{margin=2.5cm}
\usepackage{setspace}
\setstretch{1.15}

% -------------------- Colores --------------------
\usepackage{xcolor}
\definecolor{tituloazul}{RGB}{31,100,160}
\definecolor{barraazul}{RGB}{210,230,250}

% -------------------- Links clickeables --------------------
\usepackage[
colorlinks=true,
linkcolor=tituloazul,
citecolor=tituloazul,
urlcolor=tituloazul
]{hyperref}

% -------------------- Estilo de secciones --------------------
\usepackage{titlesec}
\titleformat{\section}
{\color{tituloazul}\large\bfseries}
{\thesection}{1em}{}
\titleformat{\subsection}
{\color{tituloazul}\normalsize\bfseries}
{\thesubsection}{1em}{}
\titleformat{\subsubsection}
{\color{tituloazul}\normalsize\bfseries}
{\thesubsubsection}{1em}{}

% -------------------- Encabezados automáticos --------------------
\usepackage{fancyhdr}
\pagestyle{fancy}

\fancyhf{}
\fancyhead[L]{\textcolor{tituloazul}{Electrónica de Potencia}}
\fancyhead[R]{\textcolor{tituloazul}{Práctica 1}}
\fancyfoot[C]{\thepage}

% -------------------- Figuras --------------------
\usepackage{graphicx}
\graphicspath{{figuras/}}

\usepackage{caption}
\usepackage{subcaption}
\usepackage{float}

% -------------------- Matemáticas --------------------
\usepackage{amsmath}
\usepackage[spanish,capitalize,nameinlink]{cleveref}
\crefname{figure}{figura}{figuras}
\Crefname{figure}{Figura}{Figuras}
\crefname{equation}{ecuación}{ecuaciones}
\Crefname{equation}{Ecuación}{Ecuaciones}
\crefname{table}{tabla}{tablas}
\Crefname{table}{Tabla}{Tablas}

\usepackage{siunitx}

\sisetup{
per-mode=symbol,
detect-all
}

% -------------------- Tablas profesionales --------------------
\usepackage{booktabs}
%\usepackage{tabularx}

% -------------------- Referencias IEEE --------------------
\usepackage{cite}

% -------------------- Grafica circuitos --------------------
\usepackage[spanish]{babel}
\usepackage[siunitx]{circuitikz}
\usetikzlibrary{babel} % Crucial para evitar conflictos con el idioma español
\setlength{\headheight}{15pt} % Corrige el warning de fancyhdr
\usepackage{latexsym} % para el símbolo de polaridad de la batería

\newcommand{\estudiantes}{Jose Leonardo Perez Poveda, Sergio Hoyos Mejía, Estudiante 3}
% -------------------- Documento --------------------
\begin{document}


% -------------------- Título --------------------
\begin{center}
    {\Huge \textcolor{tituloazul}{\textbf{Práctica 1}\\ Interruptores de Electrónica de Potencia y Circuitos de Disparo. Parte 1}}
    \vspace{0.6cm}

    {\textbf{\estudiantes}}
    \vspace{0.5cm}

    \textit{Ingeniería Electrónica}\\
    \textit{Universidad de Antioquia}\\
    \textit{Medellín-Colombia}
\end{center}

\vspace{0.5cm}

% -------------------- Resumen --------------------
\noindent\textcolor{tituloazul}{\textbf{Resumen:}} Esta práctica se enfoca en la caracterización y control de interruptores de estado sólido básicos en electrónica de potencia (MOSFET e IGBT). Se analizará la interacción entre señales de control de baja potencia generadas por un microcontrolador y las etapas de conmutación, mediante la implementación física de circuitos de disparo (Gate Drivers) con aislamiento galvánico (optoacopladores y optodrivers). Se validará el funcionamiento de los controladores de compuerta  y la respuesta dinámica de los interruptores mediante mediciones de laboratorio, comparando los tiempos de conmutación, caídas de tensión y el comportamiento de los dispositivos en condiciones reales de operación.


% -------------------- Objetivos --------------------
\section{Objetivo General}
Implementar circuitos de disparo para dispositivos MOSFET e IGBT, operando como interruptores, controlando la potencia suministrada a cargas eléctricas, mediante un microcontrolador.

\section{Objetivos Específicos}
    \begin{itemize}
        \item Analizar el funcionamiento de los dispositivos MOSFET e IGBT variando sus señales de conmutación e identificando sus diferencias en aplicaciones de electrónica de potencia.
        \item Implementar aislamiento galvánico entre las etapas de control y potencia mediante el uso de optoacopladores, como medida de protección, separando el control de la etapa de potencia.
        \item Controlar la potencia suministrada a las cargas eléctricas a través de la modulación por ancho de pulso (PWM).
    \end{itemize}

\newpage % salto de pagina

% -------------------- Preinforme --------------------
\section{Preinforme}
En esta sección se presenta el avance desarrollado hasta el momento, siguiendo la pauta solicitada para el preinforme.

\subsection{Consultar}

\subsubsection{Optoacoplador 6N135}
El 6N135 es un optoacoplador de un canal para alta velocidad. Internamente usa un LED de entrada y una etapa de salida con transistor NPN de colector abierto, lo que permite aislamiento galvánico entre la etapa de control y la etapa de potencia. Su principio de funcionamiento consiste en convertir la señal eléctrica de entrada en luz y luego volverla a convertir en señal eléctrica en la salida aislada.

Parámetros destacados reportados en la consulta:
\begin{itemize}
    \item Corriente directa del LED: \SI{25}{mA}.
    \item Corriente pico directa del LED: \SI{50}{mA}.
    \item Voltaje inverso del LED: \SI{5}{V}.
    \item Corriente promedio de salida: \SI{8}{mA}.
    \item Corriente pico de salida: \SI{16}{mA}.
    \item Rango de voltaje de alimentación/salida: \SIrange{-0.5}{15}{V}.
    \item Voltaje de aislamiento: \SIrange{2500}{5000}{V_{rms}}.
\end{itemize}

\subsubsection{Optodriver TLP350}
El TLP350 es un optoacoplador con etapa de potencia para manejo de compuerta en MOSFET e IGBT. A diferencia del 6N135, su salida está diseñada para entregar y absorber corriente de compuerta, facilitando encendido y apagado rápidos del interruptor.

Parámetros destacados reportados en la consulta:
\begin{itemize}
    \item Corriente de entrada del LED: \SI{20}{mA}.
    \item Corriente pico de entrada: \SI{1}{A}.
    \item Voltaje inverso de entrada: \SI{5}{V}.
    \item Corriente pico de salida en alto: \SI{2.5}{A}.
    \item Corriente pico de salida en bajo: \SI{-2.5}{A}.
    \item Tensión máxima de alimentación: \SI{35}{V}.
    \item Voltaje de aislamiento: \SI{3750}{V_{rms}}.
\end{itemize}

\subsubsection{Circuito integrado CD40106}
El CD40106 es un integrado CMOS con seis inversores Schmitt Trigger independientes. Se usa para invertir señales lógicas y mejorar su inmunidad al ruido, gracias al efecto de histéresis en los umbrales de conmutación.

Parámetros destacados reportados en la consulta:
\begin{itemize}
    \item Rango de voltaje de entrada: \SIrange{-0.5}{18}{V}.
    \item Corriente de entrada típica muy baja (orden de \si{\micro\ampere}).
    \item Corriente de salida LOW: \SI{2.4}{mA}.
    \item Corriente de salida HIGH: \SI{-2.4}{mA}.
\end{itemize}

\subsubsection{MOSFET IRFZ44N e IGBT GP10N60}
Este apartado se encuentra \textbf{pendiente} de completar en el siguiente avance, incluyendo condiciones de operación como interruptor (región de conducción, parámetros de compuerta y limitaciones térmicas/eléctricas).

\subsection{Código PWM}
Este apartado se encuentra \textbf{pendiente}. En la siguiente versión se incorporará el código para STM32/ESP32/Raspberry Pi Pico con control serial de frecuencia (\SIrange{60}{40000}{Hz}) y ciclo útil.

\subsection{Simulaciones}
Este apartado se encuentra \textbf{pendiente}. Se incluirán simulaciones en LTspice para los circuitos de la \cref{fOptoMosfet} y la \cref{fOptodIGBT}, junto con el análisis comparativo de formas de onda.

\subsection{Fuentes consultadas}
\begin{itemize}
    \item Hoja de datos de 6N135.
    \item Hoja de datos de TLP350.
    \item Hoja de datos de CD40106.
    \item Material de apoyo sobre LED AlGaAs (ScienceDirect).
\end{itemize}

%--------------------------Práctica-----------------------

\section{Práctica}
Durante la práctica se deben implementar los circuitos de la \cref{fOptoMosfet} y la \cref{fOptodIGBT} en protoboard. Se debe ajustar la fuente de potencia a máximo 12Vdc y 5Amp. Las preguntas planteadas en esta sección se deben responder durante la práctica.

\subsection{Optoacoplador y Mosfet}
Implementar el circuito de la \cref{fOptoMosfet}, la frecuencia y ancho de pulso (duty) se le indicará durante la práctica. 


\begin{figure}[H]
\centering
\includegraphics[width=1\textwidth]{fOptoMosfet}
\caption{Circuito de disparo para Mosfet con optoacoplador de alta velocidad.}
\label{fOptoMosfet}
\end{figure}



\begin{itemize}
    \item Para la frecuencia y duty dados, obtenga las gráficas de los PWM en el microcontrolador y después del optoacoplador. ¿tienen la misma forma de onda? ¿por qué?
    \item Mida el PWM después del CD40106, ¿para qué sirve este componente en el circuito? ¿qué nombre recibe esta función?
    \item Mida el PWM en el gate-source del mosfet y comparelo con los PWM de las etapas anteriores, ¿qué es y para qué sirve un driver de disparo?
    \item ¿como se llama el arreglo formado por Q2 y Q3? ¿para qué sirve?
\end{itemize}


\subsection{OptoDriver e IGBT}
Implementar el circuito de la \cref{fOptodIGBT}, la frecuencia y ancho de pulso se le indicará durante la práctica.


\begin{figure}[H]
\centering
\includegraphics[width=1\textwidth]{fOptodIGBT}
\caption{Circuito de disparo con optodriver de alta velocidad: MOSFET vs IGBT.}
\label{fOptodIGBT}
\end{figure}

\begin{itemize}
    \item ¿Porque en este circuito no se usa el arreglo de Q2 y Q3 del primer circuito?
    \item ¿cuál es la función de R2 y R4 en el circuito?
    \item Fije el duty al 80\% y haga un barrido en frecuencia para f=60Hz, f=500Hz, f=10kHz y f=30kHz. Mida/Analice la forma del PWM directamente en la carga y mida la temperatura del MOSFET, luego repita el procedimiento con el IGBT y compare los resultados. Según su criterio, ¿cual de de los dos interruptores es mejor?
\end{itemize}



% ------------------------- Resultados -----------------

\section{Informe (Resultados y Análisis)}
Se deben presentar las gráficas/fotos y tablas obtenidas de forma organizada y técnica. El análisis no debe ser una descripción narrativa de los datos, sino una interpretación basada en fundamentos obtenidos de la experimentación con su respectiva justificación. Se debe:

\begin{itemize}
    \item Comparar los resultados obtenidos entre las diversas configuraciones o puntos de operación evaluados.
    \item Justificar el comportamiento de las variables críticas basándose en los fundamentos teóricos del sistema.
\end{itemize}

Se deben responder las preguntas planteadas para resolver durante la práctica, junto con las siguientes preguntas para el informe:

\begin{itemize}
    \item ¿Cuál es la justificación técnica para implementar aislamiento galvánico entre las etapas de control y potencia?
    \item ¿por qué hay dos tierras diferentes en los circuitos?
    \item ¿Qué son los circuitos lógicos digitales TTL?
    \item ¿Qué son los circuitos lógicos digitales CMOS?
\end{itemize}


% ------------------------- Conclusiones ---------------
\section{Conclusiones}
Deben ser deducciones técnicas basadas en los resultados y el análisis previo. Relaciónelas directamente con los objetivos de la práctica y evite frases genéricas como "se aprendió mucho", "Se aprendió a utilizar el software", "Los resultados obtenidos son iguales a la teoría", "La práctica fue muy interesante", etc.

% -------------------- Bibliografía --------------------
\bibliographystyle{IEEEtran}
\bibliography{bibliografia}

% ---------------------componentes----------------------
\newpage

\section*{Listado de Equipos, Herramientas y Componentes}

{\textbf{\estudiantes}}

\begin{table}[H]
    \centering
    \caption{Elementos a solicitar al laboratorista para la práctica.}
    \resizebox{0.8\textwidth}{!}{
    \begin{tabular}{c p{11cm} c} 
    \toprule
    \toprule
    \textit{Item} & \textit{Elemento} & \textit{Cantidad} \\
    \midrule
    1  & Osciloscopio + 1 punta de voltaje & 1 \\
    \midrule
    2  & Multímetro + 2 puntas de banana-caimán & 1\\
    \midrule
    3  & Fuente DC + 4 puntas banana-caimán & 1\\
    \midrule
    4  & Protoboard & 1\\
    \midrule
    5  & cámara térmica & 1\\
    \midrule
    6  & Pelacables & 1\\
    \midrule
    7  & Microcontrolador (STM32 o ESP32 o Raspberry PI Pico)  & 1 \\
    \midrule
    8  & Optoacoplador 6N135 & 1\\
    \midrule
    9  & Optodriver TLP350 & 1\\
    \midrule
    10  & MOSFET IRFZ44N & 1\\
    \midrule
    11  & IGBT GP10N60 & 1\\
    \midrule
    12  & Circuito integrado CD40106 & 1\\
    \midrule
    13 & Transistor 2N2222 & 2\\
    \midrule
    14 & Transistor 2N2907 & 1\\
    \midrule
    15 & Alambre UTP & 10\\
    \midrule
    16  & Bombillo 12v & 2\\
    \midrule
    17  & Resistencia 10k$\Omega$ & 2\\
    \midrule
    18  & Resistencia 220$\Omega$ & 1\\
    \midrule
    19  & Resistencia 10$\Omega$ & 1\\
    \midrule
    20  & Resistencia 4.7k$\Omega$ & 1\\
    \midrule
    21  & Resistencia 6.8k$\Omega$ & 1\\
    \midrule
    22  & Capacitor 330$p$f  & 1\\
    \midrule
    23  &   & \\
    \midrule
    24  &   & \\
    \midrule
    25  &   & \\
    \midrule
    26  &   & \\
    \midrule
    27  &   & \\
    \midrule
    \bottomrule
    \end{tabular}
    }
    \label{tElementos}
\end{table}

\end{document}